\documentclass[12pt]{article}

%\usepackage[cm]{fullpage}
\setlength{\topmargin}{-1in}
\setlength{\oddsidemargin}{-0.25in}
\setlength{\evensidemargin}{-0.25in}
\setlength{\textwidth}{7in}
\setlength{\textheight}{10in}

\usepackage{array}
\usepackage{enumitem}
\usepackage{amssymb}
\usepackage{amsmath}
\usepackage{epstopdf}
\usepackage[T1]{fontenc}
\usepackage{graphics}
\usepackage{graphicx}
\usepackage{hyperref}
\usepackage{latexsym}
\usepackage{mathrsfs}
\usepackage{multicol}
\usepackage{setspace}
\usepackage{subfigure}
\usepackage{time}
\usepackage{times}

\newcommand{\To}{\Rightarrow}
\newcommand{\del}{\nabla}
\newcommand{\R}{\mathbb{R}}
\newcommand{\C}{\mathcal{C}}
\newcommand{\N}{\mathcal{N}}
\newcommand{\A}{\mathcal{A}}
\renewcommand{\O}{\mathcal{O}}
% \renewcommand{\N}{\mathbb{N}}
\newcommand{\ben}{\begin{enumerate}}
\newcommand{\een}{\end{enumerate}}
\newcommand{\eps}{\varepsilon}
\renewcommand{\epsilon}{\varepsilon}
\newcommand{\p}{\partial}
\newcommand{\<}{\langle}
\renewcommand{\>}{\rangle}
\renewcommand{\subset}{\subseteq}
\newcommand{\cont}{\Rightarrow \Leftarrow}
\newcommand{\back}{\backslash}
\newcommand{\norm}[1]{\|{#1}\|}
\newcommand{\abs}[1]{|{#1}|}
\newcommand{\ip}[1]{\langle{#1}\rangle}
\newcommand{\m}[1]{{\mathbf{#1}}}
\newcommand{\sts}{\m{S}^\textrm{T}\m{S}}
\newcommand{\rank}{\textrm{rank}}
\newcommand{\X}{\mathcal{X}}
\newcommand{\Y}{\mathcal{Y}}
\def\ds{\displaystyle}
\def\d{\partial}
\renewcommand{\arraystretch}{1.5}
\renewcommand{\L}{\mathscr{L}}

%\includegraphics[scale = 0.65, trim = 1in 3.25in 1in 3in]{blank_axes.pdf}


\pagestyle{empty}
\begin{document}
\noindent\rule{7in}{0.4pt}\\
\begin{center}
{\sc Differential Equations\\
General Solution Versus Particular Solution}
\end{center}
\noindent\rule{7in}{0.4pt}
%%%%%%%%%%%%%%%%%%%%%%%%%%%%%%%%%%%%%%%%%%%%%%%%
\section*{General Solution vs. Particular Solution}
\doublespacing

\noindent  Sometimes, in addition to the DE itself (an equation that contains an unknown function), 
we actually have some additional information about the solution (i.e., the unknown function) itself.  
Here we discuss how to use this additional information and what it means.\\

%Maybe a different example?
\noindent  If we are studying the spread of a disease, for example, we may want to know
how many people will be infected at some point in the future given that we know how many people
are infected today.  In this case, we might have some information about the rate at which the
disease spreads, and this can help write a differential equation, because {\it rate} is associated with 
the derivative.  The number of people infected today is extra information about the solution.  If
%$I(t)$ represents the number of people infected at time $t$
\line(1,0){300} \line(1,0){300}, and there are 12 people infected today,
which we might think of as time %$t = 0,$
\line(1,0){75}, then we could write the extra information as %$I(0) = 12.$
\line(1,0){75}.  Extra information about the solution that tells us about the solution at one 
a single value of the independent variable is called an %{\bf initial condition.}\\
\line(1,0){200}.\\

\noindent  {\bf Example 4:}  Consider the differential equation $\ds \frac{dy}{dx} = \frac{x^2}{y^2}.$ 
	Suppose we happen to know that the general solution is $y = \sqrt[3]{x^3 + C}$.
	Among the family of solutions, find the solution that satisfies the additional 
	information that $y(0)=2$.\\
	\newpage
%Among the infinite number of solutions that
%satisfy the DE, we have found the one, unique solution that also satisfies the initial condition.\\


\noindent  Keep in mind that ordinarily
we would be begin by using a solution technique to find the general solution to a DE.  Since we 
have not learned any solution techniques yet, we are just practicing finding a particular solution 
once we have the general solution.\\ %Does this make sense to say after then next example, instead??

\noindent  {\bf Example 5:}  Solve
	\[ y'' - y' - 12y = 0, \hspace{0.25cm} y(0) = 4, \hspace{0.25cm} y'(0) = -2. \]
	
%\noindent  As in the previous example, we'll use the general solution we already verified in Example 1,
%$ y = c_1e^{-3t} + c_2e^{4t}, $ as our starting point.  Notice that the second piece of extra information
%gives us information about the value of the derivative, $y'$; therefore we would need to find the first derivative,
%$y' = -3c_1e^{-3t} + 4c_2e^{4t},$ if we haven't already done so.
%
%Now let's see what happens when we use the initial conditions:
%
%\begin{align*}
%y(0)	&= 4		&\Rightarrow \hspace{0.2cm} 4		&= c_1e^{-3\cdot 0} + c_2e^{4\cdot 0}\\
%	&		&						4	&= c_1 + c_2\\
%	&		&							&\\
%y'(0)	&= -5	&\Rightarrow \hspace{0.2cm} - 5	&=-3c_1e^{-3\cdot} + 4c_2e^{4\cdot}\\
%	&		&						-5	&= -3c_1 + 4c_2\\
%\end{align*}
%
%\noindent  Notice that we now have two linear equations in terms of two unknowns, $c_1$ and $c_2.$
%There are multiple ways to solve this system of equations, and the reader is encouraged to verify the
%details that show that $c_1 = 3$ and $c_2 = 1$.  %OK to omit details here?
%%We will solve for $c_2$ in the first equation
%%and then substitute into the second.
%%	\begin{align*}
%%			4	&= c_1 + c_2\\
%%	\Rightarrow c_2	&=4 - c_1\\
%%				&\\
%%			-5	&= -3c_1 + 4c_2\\
%%	\Rightarrow -5	&= -3c_1 + 4(4 - c_1)\\
%%			-5	&= -3c_1 + 16 - 4c_1\\
%%			-21	&= -7c_1\\
%%			3	&= c_1\\
%%				&\\
%%			c_2	&=4 - c_1\\
%%				&= 4 - 3\\
%%				&= 1\\
%%	\end{align*}
%\noindent  Thus, the solution we that satisfies the DE and both of the provided initial conditions is
%	\[ y = 3e^{-3t} + e^{4t}. \]

\newpage
	
\noindent  You've now seen a couple of examples that demonstrate how we can use extra information
to narrow exactly one solution from the infinite number we have in the general solution.  When we do
find exactly one solution, we call it a %{\bf particular solution}
\line(1,0){175}.  This is what we found in the previous
two examples when we applied the given initial conditions.  In fact, this same question will appear
so frequently that is has a name.
	\begin{quote}  {\it Definition:}  When we have a DE and we are given initial conditions to go along with the DE,
	we have an %{\bf initial value problem,} or an IVP.
	\line(1,0){300}.
	\end{quote}
	
\noindent  When we were finding the general solution to a DE, there were multiple ways that question
might be asked.  The same is true for the particular solution.  The following are equivalent ways to ask 
the same question.
\begin{itemize}
\item  Solve $ y'' - y' - 12y = 0, \hspace{0.25cm} y(0) = 4, \hspace{0.25cm} y'(0) = -2. $
\item  Solve  the initial value problem $ y'' - y' - 12y = 0, \hspace{0.25cm} y(0) = 4, \hspace{0.25cm} y'(0) = -2. $
\item  Find the particular solution to $ y'' - y' - 12y = 0, \hspace{0.25cm} y(0) = 4, \hspace{0.25cm} y'(0) = -2. $
\end{itemize}
	
\noindent  Keep in mind that ordinarily
we would be begin by using a solution technique to find the general solution.  Since we 
haven't learned any solution techniques yet, we are just practicing find a particular solution 
once we have the general solution.)\\

\noindent  Initial conditions are not the only kind of extra information we might have.  Sometimes
we are given information about the solution at different values of the independent variable.  These
are called boundary values and lead to boundary value problems.  We will revisit this idea in a
future chapter.\\

%%\noindent  Summary:\\
%%\begin{itemize}
%%\item  Know what a DE is
%%\item  Know what it means to be a solution to a DE
%%\ite  Be to determine whether a function is/is not a 
%%	solution to a DE
%%\item  Be able to identify independent and dependent variables in a DE
%%\item  Be able to explain the difference between the general solution and a particular solution.
%%\item  Be able to recognize various ways you might be asked to solve for the general solution 
%%and/or the particular solution
%%\item  Be able to determine the order of a DE
%%\item  Be able to determine whether a DE is linear or nonlinear
%%\end{itemize}


\end{document}